\section{Why the DGF is the most likely first candidate for agentic automation}
\label{sec:candidate}

Alongside the law, we state a separate comparative hypothesis: \textbf{among the expert control
functions of an enterprise, the DGF will be the first to be completely executed by agents.} The law
does not depend on it; the DGF could close after another function and the law would still hold. The
hypothesis rests on seven properties that few other functions combine.

\textbf{Its objects are already digital.} A DGF reviews software systems, configurations, contracts,
evidence, and institutional commitments. Repositories, identity systems, cloud control planes, ticketing
tools, and APIs already supply most of the observations and actions a gate needs. Nothing has to be lifted, driven, or examined at a bedside: every role reads documents and returns an
artifact (Table~\ref{tab:roles}).

\textbf{It is already a paved road} (Section~\ref{sec:dgf}). The representation work that blocks
automation elsewhere has been done here for human reasons.

\textbf{Its decisions concern systems and projects.} An architecture approval is not a hiring, credit,
or medical decision about a person. This does not exempt it from law (Section~\ref{sec:mandate}), but it
places the DGF further from the strictest regimes on automated decisions than most candidate functions.

\textbf{It is a cost center and a bottleneck.} Review queues delay every project of the enterprise. An internal review is not
a product feature: customers buy the firm's products, not the identity of whoever performed its
internal reviews. Where a contract or a regulator does require a named human reviewer, that is an
institutional constraint of the kind discussed in Section~\ref{sec:mandate}, not a commercial
preference. Management therefore faces both a speed incentive and a cost incentive.

\textbf{Its work is specialist work, and much of it is repetition.} The same facts are extracted from the
same documents by Procurement, Legal, Security, and IT production. How many visits are standard is an empirical frequency, measurable on historical logs. Hypothesis~H3
makes a separate prediction: the residual holds a larger share of the hours than of the cases
(Section~\ref{sec:selection}), whether or not standard cases are the majority.

\textbf{Its outputs are checkable downstream.} A condition such as ``EDR before interconnection'' or
``log export in the contract'' can be verified by the next gate or by the production system itself. A
function whose outputs can be tested is a function whose automated executor can be qualified.

\textbf{Its neighbors are already software.} Bounded governance decisions around the DGF already execute
without a fresh human judgment (Table~\ref{tab:precedents}).

The comparison class matters. Transactional back-office functions, such as invoice processing, order
creation, post-trade matching, and first-level service desks, are already largely executed by software;
they are the precedents of Section~\ref{sec:precedents}, not rivals for this hypothesis. The rivals are
the other expert control functions staffed by senior specialists: legal review, financial control and
audit, credit and underwriting decisions, human-resources decisions, clinical review. The criterion is
execution closure as defined in Section~\ref{sec:conditions}, reached for every gate of the function at
a participating site. Human-resources, credit, and clinical decisions concern persons and meet the
strictest legal regimes; legal work carries privilege and professional liability; audit is bound by
independence rules. The DGF combines high structure, digital evidence, internal impact, and management
pressure. The hypothesis fails if another expert control function reaches execution closure first. On
this hypothesis, the populations of Table~\ref{tab:populations} are the first specialist workforce
exposed.

\section{Why believe complete automation will extend to every gate?}
\label{sec:premises}

The argument is inductive and compositional, rather than a deduction from notation.
Table~\ref{tab:premises} distinguishes evidence already available from the extension that the thesis
needs.

\begin{table}[!htb]
\centering\small
\caption{Five premises, with their evidential boundary. The final column contains the substantive
burden of the complete-automation thesis.}
\label{tab:premises}
\begin{tabularx}{\linewidth}{@{}P{3.0cm}LL@{}}
\toprule
Premise & Evidence or mechanism & Remaining claim to test\\
\midrule
Contract, not occupation & Each DGF specifies evidence, criteria, outputs, and the order chosen for its
routes. & Contracts
capture material inquiry and judgment without concealing a necessary human.\\
Capability reaches gate scale & Measured software-task horizons have increased substantially. &
Comparable progress transfers to messy gate work and the required reliability tail.\\
Gates already become software & CI quality rules, cloud policy, standard changes, and transaction
processing execute bounded controls. & Automation extends from these controls to the full judgment and
exception loop.\\
Authority becomes a mandate & Authorized rules can release changes or orders without another per-case
decision. & Required commitments can be validly delegated on each DGF route.\\
Provision scales & Reusable services can spread common engineering and assurance costs. & Local
adaptation and risk costs do not preserve an economic or human-support floor.\\
\bottomrule
\end{tabularx}
\end{table}

\subsection{Capability is moving toward longer coherent work}
\label{sec:capability}

\citet{kwa2025} introduce a task-completion horizon measured in human expert time, not in the agent's
runtime. Their 2019--2025 results show an approximately seven-month doubling of the 50\% horizon on
software-related tasks, and a similar trend for the 80\% horizon at roughly one fifth of the length.
They associate the increase with better reliability, error recovery, reasoning, and tool use. This is
the strongest measured trend supporting the move from individual instructions to extended work
packages.

The later METR dashboard, updated on 8 May 2026, reports for its best-performing model a 50\% horizon
of at least sixteen hours, with a 95\% interval of 8.5--55 hours, and warns that measurements above
sixteen hours exceed the reliable range of its task suite \citep{metr2026a}. The reported point is not
a guarantee that every sixteen-hour activity can be completed.

The inference to gates requires care. A twelve-, twenty-four-, or thirty-six-hour dossier can contain
parallel work by several specialists; it is not automatically one coherent task of that duration. The
12--36-hour examples below are synthetic reference workloads, not measured DGF horizons. METR's tasks
are mainly self-contained software engineering, machine-learning, and cybersecurity problems; its
documentation emphasizes domain variation over orders of magnitude, low-context human baselines, worse
transfer to messy work, and error bars of roughly a factor of two \citep{metr2026a,metr2026b}. It does
not estimate 99\% horizons.

These qualifications change how the evidence is used, not its relevance. A trend from short tasks to
multi-hour, tool-using work is evidence against a fixed boundary at a document summary. The thesis is
that the same combination of improved reasoning, access, recovery, and verification closes the larger
contracts. Sections~\ref{sec:milestones} and~\ref{sec:testing} expose the additional assumptions and
tests.

\subsection{The inductive precedent: some decisions already run as software}
\label{sec:precedents}

Table~\ref{tab:precedents} documents operational mechanisms in the same organizational neighborhood as
the DGF. The important observation is not a vendor's adoption claim. It is that an input can already be
checked against an authorized condition and cause an effective action without a fresh human judgment
in that bounded case.

\begin{table}[!htb]
\centering\small
\caption{Existing automation near the gate boundary. Documentation supports the stated mechanism,
without implying zero supporting labor or full automation of the wider function.}
\label{tab:precedents}
\begin{tabularx}{\linewidth}{@{}P{3.3cm}LL@{}}
\toprule
Primary example & Software-executed transformation & Boundary of the precedent\\
\midrule
CI/CD quality gates \citep{sonar2026} & Code-analysis measurements and conditions produce pass/fail;
pipeline or merge policy can enforce the result. & Encoded quality criteria, not every design or
security judgment.\\
Cloud policy-as-code \citep{mspolicy2025} & A resource request matching a deny policy is rejected
before resource-provider execution. & Enforces a rule; does not establish that the rule covers all
risks.\\
Standard changes \citep{atlassian2026} & A repeatable, documented, low-risk change can be
pre-authorized rather than approved anew. & Preauthorization supplies the mandate; automated execution
still needs an implementation.\\
Continuous deployment \citep{awscd2026} & A qualifying change proceeds automatically into production,
unlike delivery awaiting explicit release approval. & Tests, configuration, rollback, and supporting
operations remain relevant.\\
Touchless purchasing \citep{oracle2026} & An approved requisition can become a purchase order and be
communicated without a procurement agent's intervention. & Requisition approval and agreement formation
are prior inputs, not shown to be autonomous.\\
Compliance evidence \citep{awsaudit2026} & Service data are collected and associated with control
evidence. & The service explicitly does not assess compliance itself.\\
Post-trade processing \citep{dtcc2026} & Matching enables automatic trade affirmation and downstream
processing. & A bounded post-trade path, not the entire back office or all exceptions.\\
Automated underwriting \citep{fanniemae2026} & Application and credit information produce a risk
assessment and recommendation. & The lender retains underwriting judgment and responsibility.\\
\bottomrule
\end{tabularx}
\end{table}

These precedents refute the proposition that a governance-related decision necessarily requires an
employee to handle every instance. Some implement the decision itself; others automate evidence or
recommend an outcome while retaining a human. That contrast is useful: it identifies which part of the
contract remains open instead of labelling every use of software a completed gate.

The next step is therefore an extension, not a new direction of history. Deterministic controls
already translate accepted requirements into executable tests. Agents add interpretation, evidence
pursuit, exception diagnosis, and revision of a proposed plan. The favorable induction is that a
growing composition of such operations replaces the whole gate. Its strongest rival is that
unstructured judgment or assurance remains an irreducible last task. Section~\ref{sec:objections} tests
that rival rather than assuming it away.

\subsection{From an individual signature to a standing mandate}
\label{sec:mandate}

Authorization is not the same as a handwriting event. Continuous deployment and touchless order
creation illustrate execution under rules and prior approvals \citep{awscd2026,oracle2026}. A DGF can
similarly move an eligible class of decisions from repeated individual approval to a standing
delegation. The delegation specifies case class, allowed actions, budget and risk limits, evidence
requirements, expiry, suspension, and the principal on whose behalf the system acts.

At the general gate, the go decision and the sponsors' commitment are then a recorded, valid exercise of
that mandate. The gate remains and the commitment is effective; no person is impersonated and no machine declares itself a corporate officer. Cases
outside the mandate remain unresolved or are refused under the contract. A mandate with a required
human exception path has only partially automated the gate; its extension to every case class is part
of the thesis.

The legal limit is substantive. NIS2 Article~20 requires management bodies to approve and oversee
cybersecurity risk-management measures and addresses their liability and training. DORA Article~5
assigns management bodies responsibility for the ICT risk-management framework; Article~6(10) permits
outsourcing verification tasks while retaining the financial entity's responsibility
\citep{nis2,dora}. These texts support distinguishing execution from responsibility, but do not
authorize eliminating all human governance. Required continuing board work belongs in $B_A$. Under a
regime that retains it, a zero-support claim for that broad boundary fails unless the requirement
itself validly changes.

Nor is ``a decision about systems'' a general exemption. GDPR Article~22 concerns solely automated
decisions about people with legal or similarly significant effects, subject to specified exceptions and
safeguards; ordinary architecture approval is not automatically such a decision, but a DGF service can
affect access, employment, or other protected interests \citep{gdpr}. The AI Act's high-risk
classification and human-oversight provisions also depend on intended use and applicable conditions,
not the label DGF \citep[Arts.~6, 14 and Annex~III]{aiact}. Jurisdiction, national implementation,
effective dates, and sector rules must be assessed for the actual service. No legal opinion or
universal permission is asserted here.

The argument for migration is thus operational: the authorized principal can define a class-level rule
and let software execute it. The argument against complete support closure is institutional: some
continuing human obligations may survive that migration. Both are observable. Maintaining this
distinction strengthens the automation thesis by identifying what must change, rather than treating
liability as either a magical human advantage or an irrelevant detail.

\subsection{Shared provision can make a low-volume route economical}
\label{sec:provision}

A five-visit Integrate route may be uneconomic when it must build and maintain its own agent
platform. That does not imply that it remains uneconomic when buying a common service. Let $F_s$ be
fixed, shareable nonlabor platform cost, $N$ comparable customers, $F_c$ customer-specific fixed cost,
$a$ variable nonlabor cost per visit, and $\lambda_c$ visits at customer $c$. Under an equal allocation
convention,
\begin{equation}
\bar c_{c,\mathrm{platform}} = a + \frac{F_s/N + F_c}{\lambda_c}. \label{eq:platformcost}
\end{equation}
The common term falls with scale, while local integration does not. Falling provider cost reaches the
customer only if pricing passes through some of the saving. The same distinction applies to human
upkeep. Allocating shared hours $B_s$ and customer-specific hours $B_c$ gives
\begin{equation}
b_{A,c} = \frac{B_s/N + B_c}{\lambda_c h_{0,c}}. \label{eq:sharedupkeep}
\end{equation}

\begin{table}[!htb]
\centering\small
\caption{Synthetic low-volume provision: 60 common human upkeep hours per month plus six
customer-specific hours, five visits per customer, and 24 baseline hours per visit. Provider labor has
not disappeared; it is allocated explicitly.}
\label{tab:provision}
\begin{tabular}{@{}r r r r r@{}}
\toprule
Customers & Support hours/month/customer & Support hours/visit & Total hours/visit & $\rho$\\
\midrule
\rowsinput{generated/table_provision_rows}
\bottomrule
\end{tabular}
\end{table}

Table~\ref{tab:provision} applies Equation~\eqref{eq:sharedupkeep} to the five-visit Integrate route of
Section~\ref{sec:routes}, whose direct residual effort is 21.62 hours per visit. Alone, the customer
faces a workload ratio of 1.451; sharing the sixty common hours across 1,000 customers while keeping
all six local hours lowers it to 0.951. The verdict on a low-volume route therefore turns on
whether it must carry its own policies, tests, and evaluations or can share them. The table holds the
sixty common hours constant as customers are added, which is the favorable case: the empirical
question is whether shared support hours grow less than proportionally with the number of customers.

Amortization lowers a buyer's allocated burden; it does not eliminate the provider's workers. Across
all customers the shared hours still sum to $B_s$, and the local hours to $\sum_c B_c$. Full support
closure requires automating or removing the required work itself. Economies of scale make the
investment plausible; they are not a shortcut to zero global labor.
